\documentclass[12pt]{article}
\usepackage{notes-project}
\usepackage{braket} 

\begin{document}

\begin{center}
  {\LARGE\bfseries Chapter 3: Formalism}\\[0.25em]
  {\large Quantum Mechanics}\\[0.15em]
\end{center}

% Key definitions from Griffiths Ch. 3 typed as needed for cross-references.

\section{Commutators}

\begin{dfn}[Commutator]{commutator}
  $\comm{\hat{A}}{\hat{B}} := \hat{A}\hat{B} - \hat{B}\hat{A}$.
\end{dfn}

\begin{rem}[Notation]{commutator-notation}
  Others usually write commutators as $[\hat{A}, \hat{B}]$.
\end{rem}

\begin{thm}[Commuting Operators Share Eigenvectors]{commuting-share-eigenvectors}
  If $\comm{\hat{A}}{\hat{B}} = 0$, then $\hat{A}$ and $\hat{B}$ share a common eigenbasis.
\end{thm}

\begin{proof}
  {\color{blue}\textbf{TODO}}
\end{proof}

\begin{thm}[Linearity of the Commutator]{commutator-linearity}
  For any operators $\hat{A}, \hat{B}, \hat{C}$,
  \[
    \comm{\hat{A} + \hat{B}}{\hat{C}} = \comm{\hat{A}}{\hat{C}} + \comm{\hat{B}}{\hat{C}}.
  \]
\end{thm}

\begin{proof}
  {\color{blue}\textbf{TODO}}
\end{proof}

\begin{thm}[Leibniz Rule for the Commutator]{commutator-leibniz}
  For any operators $\hat{A}, \hat{B}, \hat{C}$,
  \[
    \comm{\hat{A}\hat{B}}{\hat{C}} = \hat{A}\,\comm{\hat{B}}{\hat{C}} + \comm{\hat{A}}{\hat{C}}\,\hat{B}.
  \]
\end{thm}

\begin{proof}
  {\color{blue}\textbf{TODO}}
\end{proof}

\begin{thm}[Canonical Commutation Relation]{canonical-commutation}
  For the position and momentum operators,
  \[
    \comm{\hat{x}}{\hat{p}} = i\hbar.
  \]
\end{thm}

\begin{proof}
  {\color{blue}\textbf{TODO}}
\end{proof}

\begin{thm}[Hamiltonian Commutators]{hamiltonian-commutators}
  For a 1D Hamiltonian $\hat{H} = \dfrac{\hat{p}^2}{2m} + V(\hat{x})$ (\xref{ch02::def:hamiltonian-1d}),
  \[
    \comm{\hat{H}}{\hat{x}} = -\frac{i\hbar}{m}\,\hat{p}, \qquad
    \comm{\hat{H}}{\hat{p}} = i\hbar\, V'(\hat{x}).
  \]
\end{thm}

\begin{proof}
  {\color{blue}\textbf{TODO}}
\end{proof}

\section{Hermitian Operators}

\begin{dfn}[Hermition Conjugate] 
  The \emph{Hermitian conjugate} of an operator $\hat{O}$ is denoted $\hat{O}^{\dagger}$ and is defined by the property
  \[\forall f, g, \braket{f|\hat{O}g} = \braket{\hat{O}^{\dagger}f|g}\]
\end{dfn} 

\begin{rem}[Hermitian Conjugate as Matrix Transpose] 
  If $\hat{O}$ is represented as a matrix, then the Hermitian conjugate is the complex conjugate transpose of that matrix. ie, 
  \[H^\dagger = \left(H^{\text{T}}\right)^{*}\]
\end{rem}

\begin{dfn}[Hermitian Operator] 
  An operator $\hat{O}$ is \emph{Hermitian} if it is equal to its Hermitian conjugate, i.e. $\hat{O}^{\dagger} = \hat{O}$.
\end{dfn} 

\section{Unitary Operators} 

\begin{dfn}[Unitary Operator]
  An operator $\hat{U}$ is \emph{unitary} if it satisfies the property $\hat{U}^{\dagger}\hat{U} = \hat{U}\hat{U}^{\dagger} = \hat{I}$, where $\hat{I}$ is the identity operator.
\end{dfn}

\section{Conjugation and Similarity Transforms}

\begin{dfn}[Conjugation Action]{conjugation-action}
  The \emph{conjugation action} of an operator $\hat{U}$ on another operator $\hat{A}$ is defined as
  \[
    \hat{A} \;\longmapsto\; \hat{U}^{\dagger}\,\hat{A}\,\hat{U}
  \]

  After applying the operator $\hat U$ to $\ket{\psi}$, the expected value of $\hat{A}$ becomes
  \[
   \bra{\hat U \psi}\hat{A}\ket{\hat U \psi} = \bra{\psi}\hat{U}^{\dagger}\hat{A}\hat{U}\ket{\psi}.
  \]

  So the conjugation action gives the transformed operator $\hat{U}^{\dagger}\hat{A}\hat{U}$ whose expected value on the original state $\ket{\psi}$ matches $\hat{A}$'s expected value on the transformed state $\hat{U}\ket{\psi}$.
\end{dfn}

\section{Functions of Operators}

\begin{dfn}[Holomorphic Functional Calculus]{operator-function}
  For every function $f$ analytic in a neighborhood of some center $c \in \mathbb{C}$, with Taylor expansion
  \[
    f(z) \;=\; \sum_{n=0}^{\infty} \frac{f^{(n)}(c)}{n!}\,(z - c)^n,
  \]
  and every operator $\hat{A}$, define
  \[
    f(\hat{A}) \;:=\; \sum_{n=0}^{\infty} \frac{f^{(n)}(c)}{n!}\,(\hat{A} - c\,\hat{I})^n
  \]
  (subject to convergence; $\hat{I}$ is the identity operator). This blanket recipe defines a canonical map $f \longmapsto f(\hat{A})$ from analytic scalar functions to operators, known as the \emph{holomorphic functional calculus}.
\end{dfn}

\begin{ex}[Operator Exponential]{operator-exponential}
  Take $f(z) = e^z$ with center $c = 0$. Since $(e^z)^{(n)} = e^z$ and $e^0 = 1$, we have $f^{(n)}(0) = 1$ for all $n$, so the scalar Taylor series reads
  \[
    e^{z} \;=\; \sum_{n=0}^{\infty} \frac{z^n}{n!}.
  \]
  Feed this into the operator recipe. With $c = 0$, the shift $\hat{A} - c\hat{I} = \hat{A}$, giving
  \[
    e^{\hat{A}} \;:=\; \sum_{n=0}^{\infty} \frac{f^{(n)}(0)}{n!}\,\hat{A}^n \;=\; \sum_{n=0}^{\infty} \frac{\hat{A}^n}{n!}.
  \]
\end{ex}

\begin{rem}[$f(\hat{A})$ Is a Constructed Operator]{operator-function-meaning}
  $f(\hat{A})$ is not literal function application: $\hat{A}$ is an operator, not a scalar. The symbol names a \emph{new operator} built from $f$'s Taylor coefficients (scalars) and $\hat{A}$'s powers, using only operations that operators support (powers, scalar multiplication, addition).
\end{rem}

\begin{rem}[$f(\hat{A})$ Is a First-Class Operator]{operator-function-first-class}
  \begin{itemize}[nosep]
    \item Once built, $f(\hat{A})$ is an ordinary operator: add, compose, conjugate, or nest as usual. Commutes with $\hat{A}$ and any $g(\hat{A})$; may not commute with operators outside $\{\hat{A}\}$'s commutant.
    \item The choice of center $c$ is a computational convenience only --- $f(\hat{A})$ is intrinsic, and different centers give the same operator where their series converge.
  \end{itemize}
\end{rem}

\begin{thm}[Action on Eigenvectors]{operator-function-eigenvector}
  If $\hat{A}\ket{\lambda} = \lambda\ket{\lambda}$, then $f(\hat{A})\ket{\lambda} = f(\lambda)\ket{\lambda}$.
\end{thm}

\begin{proof}
  Since $\hat{A}^n\ket{\lambda} = \lambda^n\ket{\lambda}$, we have $(\hat{A} - c\hat{I})^n\ket{\lambda} = (\lambda - c)^n\ket{\lambda}$, so
  \[
    f(\hat{A})\ket{\lambda} = \sum_{n=0}^{\infty}\frac{f^{(n)}(c)}{n!}(\lambda - c)^n\ket{\lambda} = f(\lambda)\ket{\lambda}.
  \]
\end{proof}

\section{Ehrenfest's Theorem}

\begin{thm}[Generalized Ehrenfest's Theorem]{ehrenfest}
  For any observable $\hat{A}$,
  \[
    \frac{d}{dt}\langle\hat{A}\rangle = \frac{1}{i\hbar}\langle\comm{\hat{A}}{\hat{H}}\rangle + \left\langle\frac{\partial\hat{A}}{\partial t}\right\rangle.
  \]
\end{thm}

\end{document}
